% Generated draft: requires mathematical and source review.
\documentclass[11pt]{article}
\usepackage[margin=25mm]{geometry}
\usepackage{fontspec}
\setmainfont{Latin Modern Roman}
\usepackage{amsmath,amssymb,mathtools,unicode-math}
\setmathfont{Latin Modern Math}
\usepackage{enumitem,needspace,xurl,hyperref}
\setlength{\parindent}{0pt}
\setlength{\parskip}{5pt}
\setlength{\emergencystretch}{3em}
\newcommand{\N}{\mathbb{N}}
\newcommand{\Z}{\mathbb{Z}}
\newcommand{\Q}{\mathbb{Q}}
\newcommand{\R}{\mathbb{R}}
\newcommand{\C}{\mathbb{C}}
\newcommand{\F}{\mathbb{F}}
\newcommand{\A}{\mathbb{A}}
\newcommand{\PP}{\mathbb{P}}
\newcommand{\E}{\mathbb{E}}
\newcommand{\eps}{\varepsilon}
\newcommand{\dd}{\,\mathrm{d}}
\begin{document}
\section*{UnsolvedMath ID: 30001627}
\textbf{Source problem code:} catalogue entry
\subsection{Definitions and mathematical statement}
An algebraic variety is quasi-affine if it is isomorphic to an open subvariety of an affine variety over its base field. A rational algebraic \(G\)-module is a vector space equipped with an algebraic linear action of \(G\). In the usual finite-dimensional interpretation, this action is a morphism \(G\to\operatorname{GL}(V)\). For a homogeneous space equipped with an appropriate locally compact realization, unimodularity means that it admits a nonzero invariant Radon measure. The record does not specify the base field, the realization on which the measure is defined, or explicitly the dimension of \(V\); these qualifications remain unresolved.

\textbf{Supplied question.} Let \(G\) be an algebraic group and \(H\subseteq G\) an algebraic subgroup. Suppose that \(Z=G/H\) is unimodular and quasi-affine, and let \(\mu_Z\) be the \(G\)-invariant measure under consideration. Does there exist a rational \(G\)-module \(V\) and a \(G\)-equivariant embedding
\[
\iota:Z\hookrightarrow V,\qquad \iota(gz)=g\iota(z),
\]
such that \(\iota_*\mu_Z\) is a tempered distribution on \(V\)? This is the target stated by the supplied corrected formulation, subject to the unresolved choices described below.

To make the analytic requirement explicit, suppose an interpretation has been fixed in which \(Z\) carries the specified measure and \(V\) has a finite-dimensional underlying real vector space \(E\). After choosing coordinates \(E\cong\R^n\), the Schwartz space \(\mathcal S(E)\) consists of smooth functions \(\varphi\) for which
\[
p_{\alpha,\beta}(\varphi)=\sup_{x\in\R^n}\left|x^\alpha\partial^\beta\varphi(x)\right|<\infty
\quad\text{for every }\alpha,\beta\in\Z_{\geq0}^{n}.
\]
A tempered distribution is a continuous linear functional on this space with its Schwartz topology. Under this interpretation, the requested pushforward is the functional
\[
\langle\iota_*\mu_Z,\varphi\rangle
=\int_Z\varphi(\iota(z))\,\dd\mu_Z(z),
\qquad \varphi\in\mathcal S(E),
\]
with the requirement that these integrals be absolutely convergent and that the resulting functional be continuous on \(\mathcal S(E)\). This describes the usual analytic meaning; it does not supply the missing base-field or realization hypotheses.

The original wording calls the measure on \(V\) a ``pullback'' of \(\mu_Z\), whereas the corrected wording calls it a pushforward. The displayed integral follows the corrected wording. The category of embedding is also unspecified: an algebraic locally closed immersion and a smooth equivariant embedding are different requirements. No choice between them is asserted here.
\subsection{Short English statement}
The supplied question asks whether every unimodular quasi-affine homogeneous space \(G/H\) admits an equivariant embedding into a rational representation for which its invariant measure defines a tempered distribution on the representation space. The base field, analytic realization, dimension convention, and category of embedding require clarification. The supplied dated curation labels the problem partially solved; this draft does not independently verify that assessment.
\subsection{Sources}
\begin{itemize}
\item Local UnsolvedMath registry: supplied complete source record, permanent ID \(30001627\), source problem code \texttt{catalogue entry}, entitled \emph{Tempered Embeddings of Unimodular Quasi-Affine Spaces}. The project identifies \texttt{data/unsolvedmath/problems.json} as a registry and \texttt{catalog/problems.json} as an additional source of existing records; their contents were not independently inspected for this draft.
\item Supplied report citation: \emph{Representation Theory and Harmonic Analysis (2011)}, Oberwolfach Reports, DOI \texttt{10.4171/owr/2010/52}, pp.~3043--3083. Original OWR report: \url{https://doi.org/10.4171/owr/2010/52}. The record describes Oberwolfach Reports as documenting mathematical workshops and their open-problem discussions. It records proposed year \(2010\) and supplies no named proposer.
\item Supplied literature reference, labeled ``nilpotent-orbit temperedness'': \url{https://arxiv.org/abs/math/0209151}. No author names or bibliographic title for this reference are supplied in the record.
\item Literature assessment dated exactly \texttt{2026-08-21}: supplied dated curation, not independently verified. The supplied status is \texttt{partially\_solved}. The assessment says that important special cases are proved, including invariant measures on reductive adjoint orbits and nilpotent-orbit settings, while the stated general rational-module embedding conjecture for unimodular quasi-affine spaces remains unresolved. The background explicitly describes this as dated literature triage derived from the supplied OWR curation and requests independent verification of time-sensitive status and cited claims.
\item The registry labels the statement \texttt{corrected\_verified} and reports that it was checked against the official Oberwolfach report, preserving the stated hypotheses, notation, and open target. This is supplied verification metadata, not a source check performed for this draft.
\end{itemize}
\end{document}
